
\begin{frame}
	\onslide<1->{
		\frametitle{State Space Representation} 
		\framesubtitle{Example:  DC Motor with Inertial Load: An Electromechanical System}
	}
	\begin{center}
		\includegraphics[width=.9\textwidth,height=.6\textheight]{\dayonesegfour/dc_motor_w_inertial_load.png}
	\end{center}
	\onslide<2->{\center --------------------------------------------------------------------}
\end{frame}

	
\begin{frame}
	\onslide<1->{
		\frametitle{State Space Model}
		\framesubtitle{DC Motor with Inertial Load}
	}
	\onslide<2->{
	\begin{block}{System Description}
	\small	Consider a DC motor with an inertial load. The key variables and constants are as follows:
		\begin{align*}
		R & : \text{Motor resistance} \\
		L & : \text{Motor inductance} \\
		\theta & : \text{Rotational position of the load} \\
		\omega & : \text{Rate of change of }\theta \\
		K_f & : \text{Frictional constant} \\
		\tau & : \text{Torque} \\
		K_m & : \text{Armature constant}\\
		v_{\text{app}} & : \text{Applied voltage} \\
		v_{\text{emf}} & : \text{Back EMF} \\
		K_b & : \text{EMF constant}
		\end{align*}
	\end{block}
	}		
		
\end{frame}
	
	\begin{frame}
		
		\onslide<1->{
			\frametitle{How the Motor Generates Back EMF} 
			\framesubtitle{Electromagnetic Induction and Rotational Motion}  
		}
		\begin{itemize}
			\onslide<2->{
				\item Electromagnetic Induction
			}
			\onslide<3->{
				\begin{itemize}
					\item Changing magnetic field induces EMF in the windings
					\item Faraday's Law of Electromagnetic Induction
				\end{itemize}
			}
			\onslide<4->{
				\item Rotational Motion
			}
			\onslide<5->{
				\begin{itemize}
					\item Armature rotates within a magnetic field
					\item Windings cut through magnetic field lines
				\end{itemize}
			}
			\onslide<6->{
				\item Opposition to Applied Voltage
			}
			\onslide<7->{
				\begin{itemize}
					\item Back EMF acts against applied voltage
					\item Proportional to motor speed: \( v_{\text{emf}} = K_b \cdot \omega \)
				\end{itemize}
			}
		\end{itemize}
		\onslide<8->{\center --------------------------------------------------------------------}
		
	\end{frame}
	
	\begin{frame}
		
		\onslide<1->{
			\frametitle{Why the Motor Generates Back EMF} 
			\framesubtitle{Motor-Generator Duality and Energy Conservation}  
		}
		\begin{itemize}
			\onslide<2->{
				\item Motor-Generator Duality
			}
			\onslide<3->{
				\begin{itemize}
					\item Motor converts electrical energy to mechanical energy
					\item Generator converts mechanical energy to electrical energy
				\end{itemize}
			}
			\onslide<4->{
				\item Energy Conservation
			}
			\onslide<5->{
				\begin{itemize}
					\item Energy conversion processes are balanced
					\item Back EMF ensures energy conservation
				\end{itemize}
			}
			\onslide<6->{
				\item Regulation of Motor Speed
			}
			\onslide<7->{
				\begin{itemize}
					\item Back EMF provides feedback to control speed
					\item Higher speed increases back EMF, reducing net voltage
				\end{itemize}
			}
		\end{itemize}
		\onslide<8->{\center --------------------------------------------------------------------}
		
	\end{frame}
	
	\begin{frame}
		
		\onslide<1->{
			\frametitle{Back EMF Summary} 
			\framesubtitle{Back EMF Generation and its Effects}  
		}
		\begin{itemize}
			\onslide<2->{
				\item Back EMF is due to electromagnetic induction
			}
			\onslide<3->{
				\item Acts in opposition to applied voltage
			}
			\onslide<4->{
				\begin{itemize}
					\item Provides natural speed regulation
					\item Ensures energy conservation
				\end{itemize}
			}
			\onslide<5->{
				\item Motor can function as a generator
			}
			\onslide<6->{
				\begin{itemize}
					\item Converts mechanical energy back to electrical energy
					\item Dual functionality of motors and generators
				\end{itemize}
			}
		\end{itemize}
		\onslide<7->{\center --------------------------------------------------------------------}
		
	\end{frame}

	
\begin{frame}
	\onslide<1->{
	\frametitle{State Space Model}
	\framesubtitle{DC Motor with Inertial Load}
}
\onslide<2->{
		\small State space equations:
		
		\[
		\frac{d\theta}{dt} = \omega
		\]
		
		\[
		L\frac{di}{dt} = v_{\text{app}} - v_{\text{emf}} - Ri
		\]
		
		\[
		J\frac{d\omega}{dt} = \tau - K_f\omega
		\]
		
		\[
		v_{\text{emf}} = K_b\omega
		\]
}
\onslide<3->{		
	\small	Where:
		\begin{align*}
		\theta & : \text{Rotational position} \\
		\omega & : \text{Angular velocity} \\
		i & : \text{Current through the motor} \\
		J & : \text{Moment of inertia of the load} \\
		\tau & : \text{Applied torque}
		\end{align*}
}		
	\end{frame}

	
	\begin{frame}
	\onslide<1->{
	\frametitle{State Space Model}
	\framesubtitle{DC Motor with Inertial Load}
}
\onslide<2->{
	\small	In matrix form, the state space model can be expressed as:
		
		\[
		\left[\begin{array}{c}
		\frac{d\theta}{dt} \\
		\frac{di}{dt} \\
		\frac{d\omega}{dt}
		\end{array}\right]
		=
		\left[\begin{array}{ccc}
		0 & 0 & 1 \\
		0 & -\frac{R}{L} & -\frac{K_b}{L} \\
		0 & \frac{K_b}{J} & -\frac{K_f}{J}
		\end{array}\right]
		\left[\begin{array}{c}
		\theta \\
		i \\
		\omega
		\end{array}\right]
		+
		\left[\begin{array}{c}
		0 \\
		\frac{1}{L} \\
		0
		\end{array}\right]
		v_{\text{app}}
		\]
	}
\onslide<3->{
	The next slide shows the results for an applied voltage of 1.0V and various nominal parameter values.	
}
	\onslide<4->{\center --------------------------------------------------------------------}
	\end{frame}
	

\begin{frame}
	\onslide<1->{
		\frametitle{State Space Representation} 
		\framesubtitle{Example:  DC Motor with Inertial Load}
	}
	\begin{center}
		\includegraphics[width=.9\textwidth,height=.6\textheight]{\dayonesegfour/dc_motor_w_inertial_load_plot.png}
	\end{center}
	\onslide<2->{\center --------------------------------------------------------------------}
\end{frame}
